\chapter{Arthroscopic Repair}

\section*{Introduction \& Clinical Indications}
Arthroscopic repair is a sophisticated, minimally invasive surgical technique employed to diagnose and treat a wide spectrum of conditions affecting synovial joints. This procedure involves the insertion of a small, fiber-optic camera, known as an arthroscope, through tiny incisions (portals) into the joint space, providing direct, magnified visualization of intra-articular structures on a high-definition monitor. Through additional strategically placed portals, specialized surgical instruments are introduced to perform precise repairs, debridements, or reconstructions of damaged tissues. Commonly applied to major joints such as the knee, shoulder, hip, ankle, wrist, and elbow, arthroscopy effectively addresses pathologies including torn ligaments, meniscal tears, cartilage defects, loose bodies, and inflammatory conditions. The principal advantage of this approach over traditional open surgery is significantly reduced tissue trauma, which translates into less postoperative pain, smaller scars, and typically a much faster and smoother recovery for the patient.

\begin{itemize}
  \item Keyhole Surgery
  \item Minimally Invasive Joint Surgery
  \item Joint Endoscopy
  \item Specific Joint Arthroscopy (e.g., Knee Arthroscopy, Shoulder Arthroscopy)
  \item Arthroscopic Surgery
  \item Scope Surgery
\end{itemize}

The fundamental surgical principle of arthroscopic repair is to achieve definitive treatment of intra-articular pathology with minimal disruption to the surrounding healthy tissues. This is accomplished by utilizing small incisions (portals) through which a high-definition arthroscope and specialized instruments are introduced, allowing for precise visualization and manipulation within the joint capsule. The joint is continuously irrigated with sterile fluid, which distends the joint space, clears debris, and maintains a clear field of view.

The rationale behind this minimally invasive approach is multifaceted: it aims to reduce postoperative pain, minimize scar formation, decrease the risk of infection, and accelerate rehabilitation compared to traditional open surgical techniques. The technical goals include accurate diagnosis of the pathology, meticulous repair or reconstruction of damaged structures (e.g., meniscal repair, ligament reconstruction, labral reattachment), debridement of degenerative or inflamed tissue, and removal of loose bodies. Ultimately, the objective is to restore anatomical integrity, improve joint stability, alleviate pain, and preserve or enhance the long-term functional capacity of the joint, thereby improving the patient's quality of life and facilitating a quicker return to activities.

The genesis of arthroscopy can be traced back to the early 20th century, with pioneering efforts to visualize the interior of joints. Dr. Kenji Takagi, a Japanese surgeon, is widely recognized for his work starting in 1918, using a cystoscope to examine knee joints, though his contributions remained largely unknown internationally for decades. Concurrently, in the 1920s, Dr. Eugen Bircher in Switzerland also explored endoscopic examination of the knee. However, it was Dr. Masaki Watanabe, another Japanese surgeon, who truly revolutionized the field, earning him the title "father of modern arthroscopy." Beginning in the 1950s, Watanabe meticulously developed the first practical arthroscopes, refined techniques, and performed the first arthroscopic meniscectomy in 1962, publishing extensively on his innovations.

The widespread adoption and rapid evolution of arthroscopy gained significant momentum in the 1970s and 1980s. Key advancements included the development of fiber optics for superior illumination, the introduction of video cameras for magnified viewing on monitors, and the creation of a specialized array of small, precise surgical instruments. Surgeons such as Robert W. Jackson in Canada and Lanny L. Johnson in the United States were instrumental in popularizing the technique and expanding its applications beyond the knee to other complex joints like the shoulder, hip, and ankle. Continuous innovation in instrumentation, imaging technology, and surgical techniques has transformed arthroscopy from a diagnostic tool into a cornerstone of orthopedic surgery, enabling increasingly complex repairs and reconstructions that were once only possible through extensive open procedures.

Arthroscopic repair is indicated for a broad range of intra-articular pathologies that cause pain, instability, mechanical symptoms, or functional impairment, and have typically failed conservative management.

\begin{itemize}
  \item Repair of meniscal tears, including bucket-handle tears, radial tears, and flap tears, or partial meniscectomy for irreparable tears, particularly in the knee.
  \item Reconstruction of ruptured ligaments, most commonly the anterior cruciate ligament (ACL) and posterior cruciate ligament (PCL) in the knee, often following sports-related trauma.
  \item Repair of rotator cuff tears (supraspinatus, infraspinatus, subscapularis, teres minor tendons) in the shoulder, addressing both acute traumatic and chronic degenerative tears.
  \item Treatment of labral tears in the shoulder (e.g., Bankart lesions, SLAP lesions) or hip, often associated with instability or femoroacetabular impingement (FAI).
  \item Removal of symptomatic loose bodies, such as osteochondral fragments or synovial chondromatosis, that cause locking, catching, or pain within the joint.
  \item Chondroplasty or microfracture procedures for focal articular cartilage defects, aiming to stimulate fibrocartilage formation and alleviate pain.
  \item Synovectomy for chronic synovitis, inflammatory arthropathy (e.g., rheumatoid arthritis), or pigmented villonodular synovitis, to reduce inflammation and pain.
  \item Decompression for impingement syndromes, such as subacromial decompression for shoulder impingement or cam/pincer osteoplasty for femoroacetabular impingement in the hip.
  \item Stabilization procedures for recurrent patellar dislocation or instability, including medial patellofemoral ligament (MPFL) reconstruction.
\end{itemize}

Arthroscopic repair holds a prominent position as a primary surgical method for a vast array of intra-articular pathologies, largely owing to its minimally invasive nature, high efficacy, and favorable patient outcomes. For many acute injuries, such as meniscal tears, ACL ruptures, and most types of rotator cuff tears, arthroscopy is typically the first-line surgical intervention, especially when conservative treatments (rest, physical therapy, medication) have failed or when the severity of the injury necessitates immediate surgical correction to prevent further damage or chronic instability. It is also frequently utilized as a primary diagnostic tool when non-invasive imaging modalities, such as MRI, are inconclusive but clinical suspicion of significant intra-articular pathology remains high, allowing for simultaneous diagnosis and treatment.

In certain complex or chronic scenarios, arthroscopy may serve as a secondary or adjunctive procedure. For instance, it can be combined with open techniques for severe intra-articular fractures, where arthroscopy assists in fragment reduction and assessment of articular congruity before open fixation. It may also be considered as a salvage procedure following failed prior open surgeries or for addressing complications arising from previous interventions. However, for the overwhelming majority of its indications, arthroscopy is the preferred and primary surgical approach due to its proven benefits in terms of recovery, pain management, and functional restoration.

\section*{Relevant Surgical Anatomy}
Arthroscopic repair necessitates a thorough understanding of the intricate anatomy of the specific joint being addressed. Each synovial joint comprises articular cartilage covering bone ends, a joint capsule, synovial membrane producing lubricating fluid, and various ligaments, tendons, and menisci (in the knee) or labra (in the shoulder and hip) that provide stability and facilitate movement. Key anatomical considerations include the precise location and function of these stabilizing structures, such as the anterior and posterior cruciate ligaments (ACL, PCL) and menisci in the knee, or the rotator cuff tendons and glenoid labrum in the shoulder.

Neurovascular structures surrounding the joint are paramount for safe portal placement. For instance, in knee arthroscopy, awareness of the common peroneal nerve laterally and the popliteal artery posteriorly is critical. In the shoulder, the axillary nerve and brachial plexus are vulnerable. Lymphatic drainage, while not directly targeted, is important for understanding post-operative swelling. Surgical landmarks, such as McBurney's point for appendectomy, are less defined for arthroscopy, but precise portal placement relative to bony prominences (e.g., patella, acromion, femoral head) and palpable soft tissue structures is crucial to optimize visualization and instrument access while avoiding iatrogenic injury. The joint is typically distended with sterile saline to create a working space, making the intra-articular structures more accessible and visible.

\section*{Preoperative Workup \& Preparation}
Thorough preoperative workup and preparation are paramount to ensure patient safety, optimize surgical outcomes, and minimize complications. This typically begins with a comprehensive medical history and physical examination conducted by the surgeon and the anesthesia team to assess the patient's overall health status and identify any comorbidities. Preoperative laboratory tests usually include a complete blood count (CBC), electrolyte panel, renal and liver function tests, and coagulation studies (PT/INR, PTT). For patients over a certain age or with cardiac risk factors, an electrocardiogram (ECG) and sometimes a chest X-ray are performed, and cardiac or pulmonary clearance from a specialist may be required.

Imaging studies, such as X-rays, magnetic resonance imaging (MRI), or computed tomography (CT) scans, are almost always performed prior to surgery to confirm the diagnosis, delineate the extent of the pathology, and aid in surgical planning. Patients are typically instructed to fast for 6-8 hours for solids and 2 hours for clear liquids before surgery to prevent aspiration during anesthesia. Specific medications, particularly anticoagulants (e.g., warfarin, clopidogrel, novel oral anticoagulants) and non-steroidal anti-inflammatory drugs (NSAIDs), must be discontinued several days to a week prior to surgery, as directed by the surgeon and anesthesiologist, to reduce the risk of bleeding. A single dose of intravenous prophylactic antibiotics (e.g., cefazolin) is routinely administered within one hour before the incision to significantly reduce the risk of surgical site infection. Finally, informed consent is obtained, detailing the procedure, potential risks, benefits, alternatives, and expected recovery. Patients are also advised to arrange for transportation home and to prepare their living environment for post-operative recovery, including having crutches, slings, or other assistive devices readily available.

\textbf{Situations requiring postponement, optimization, or an alternative plan:}
\begin{itemize}
  \item Active joint infection (septic arthritis) in the joint designated for surgery, as arthroscopy can disseminate the infection and worsen the prognosis.
  \item Active skin infection, cellulitis, or open wounds overlying the proposed surgical portal sites, which significantly increase the risk of introducing bacteria into the joint.
  \item Severe joint stiffness or ankylosis that physically prevents adequate joint distension with fluid or the safe insertion and manipulation of the arthroscope and instruments.
  \item Uncorrectable coagulopathy or severe bleeding disorder that poses an unacceptably high risk of intra-articular hemorrhage or postoperative hematoma.
  \item Severe, unstable medical comorbidities (e.g., acute myocardial infarction, uncontrolled congestive heart failure, severe respiratory failure, uncontrolled diabetes with end-organ damage) that render the patient an unacceptably high anesthetic or surgical risk.
\end{itemize}

\textbf{Relative Contraindications \& Precautions:}
\begin{itemize}
  \item Severe osteoarthritis with extensive cartilage loss and bone-on-bone changes, where arthroscopy may offer only temporary symptomatic relief, and definitive arthroplasty (joint replacement) might be a more appropriate long-term solution.
  \item Morbid obesity, which can complicate patient positioning, make portal placement technically challenging, increase operative time, and elevate the risks of anesthesia, deep vein thrombosis, and wound complications.
  \item Significant peripheral vascular disease, which may compromise tissue perfusion, impair wound healing, and increase the risk of neurovascular complications.
  \item Poor patient compliance with post-operative rehabilitation protocols, which is a critical determinant of successful functional outcomes, particularly for reconstructive procedures.
  \item Previous surgery in the same joint that may have significantly altered the normal anatomy, created extensive scar tissue, or led to hardware that could impede arthroscopic access.
  \item Certain skin conditions (e.g., psoriasis, eczema) or known allergies to materials used in the procedure (e.g., antiseptic solutions, suture materials, implant components).
  \item Neuropathic joint (Charcot joint), where the underlying neurological deficit may lead to rapid joint destruction despite surgical intervention.
\end{itemize}

\section*{Surgical Technique \& Procedure}
The arthroscopic repair procedure typically commences with the administration of anesthesia, which can be general anesthesia, regional anesthesia (e.g., spinal or epidural block), or a combination, often supplemented with local anesthetic infiltration for enhanced postoperative pain control. The patient is then meticulously positioned on the operating table to ensure optimal access to the joint, which may involve supine, lateral decubitus, or beach chair positions, depending on the specific joint being operated on. For extremity joints, a pneumatic tourniquet is often applied to the limb and inflated to create a bloodless field, significantly improving visualization. The surgical area is then thoroughly prepped with an antiseptic solution (e.g., povidone-iodine or chlorhexidine) and draped to establish a sterile operating environment.

Small incisions, known as portals, usually 2-3 in number and typically less than 1 cm in length, are strategically made around the joint. The joint is then distended with sterile saline solution, and the arthroscope, connected to a camera and light source, is inserted through the initial portal. This allows the surgeon to conduct a comprehensive diagnostic survey of the entire joint, identifying all relevant pathology. Through additional portals, specialized instruments such as shavers, graspers, scissors, suture passers, and anchor delivery systems are introduced to perform the necessary repairs.

The key operative steps generally include:
\begin{itemize}
  \item \textbf{Diagnostic Arthroscopy:} Initial thorough inspection of all joint compartments to confirm the diagnosis and identify any concomitant pathologies.
  \item \textbf{Debridement:} Removal of damaged, inflamed, or degenerative tissue (e.g., partial meniscectomy, synovectomy, chondroplasty).
  \item \textbf{Preparation of Repair Site:} Freshening of tissue edges, microfracture of bone, or creation of bone tunnels for graft placement.
  \item \textbf{Repair/Reconstruction:} Placement of suture anchors into bone, passing sutures to repair torn ligaments or tendons (e.g., rotator cuff, labrum), or graft fixation for ligament reconstruction (e.g., ACL reconstruction using autograft or allograft).
  \item \textbf{Loose Body Removal:} Retrieval of any free-floating bone or cartilage fragments using graspers.
  \item \textbf{Irrigation and Final Inspection:} Thorough lavage of the joint to remove debris and a final check of the repair integrity and joint mechanics.
  \item \textbf{Closure:} Once the repair is complete and the joint is thoroughly irrigated, the instruments and arthroscope are removed. The portal incisions are then closed with sutures, sterile adhesive strips, or surgical glue, and a sterile dressing is applied, often with a compression bandage to minimize swelling. Drains are rarely used in routine arthroscopy.
\end{itemize}
Depending on the procedure, a brace, splint, or sling may be applied to protect the repair and support the limb.

\section*{Complications \& Risks}
While arthroscopic repair is generally considered a safe procedure, like any surgical intervention, it carries potential risks and complications. These can be broadly categorized into general surgical risks and procedure-specific risks.

\textbf{General Surgical Risks:}
\begin{itemize}
  \item \textbf{Bleeding:} Hemarthrosis (bleeding into the joint) can cause significant swelling, pain, and stiffness, potentially requiring aspiration.
  \item \textbf{Infection:} Septic arthritis (joint infection) is a serious but rare complication, while superficial wound infection at the portal sites is more common but usually manageable.
  \item \textbf{Anesthesia Complications:} These can range from minor issues like nausea, vomiting, and sore throat to more severe reactions such as allergic reactions, respiratory depression, or cardiovascular events (e.g., myocardial infarction, stroke).
  \item \textbf{Deep Vein Thrombosis (DVT):} Formation of blood clots in the deep veins of the leg, with the potential for pulmonary embolism (PE), a life-threatening condition.
  \item \textbf{Persistent Pain:} Chronic or intractable pain despite technically successful surgery, which may be due to underlying degenerative changes or neuropathic pain.
  \item \textbf{Scarring:} Although portal scars are small, keloid or hypertrophic scarring can occur in susceptible individuals.
\end{itemize}

\textbf{Procedure-Specific Risks:}
\begin{itemize}
  \item \textbf{Nerve Injury:} Damage to nearby sensory or motor nerves, leading to temporary or, rarely, permanent numbness, tingling (paresthesia), or weakness in the distribution of the affected nerve (e.g., saphenous nerve in knee, axillary nerve in shoulder).
  \item \textbf{Vascular Injury:} Damage to adjacent blood vessels, which can result in significant bleeding, hematoma formation, or, in severe cases, compromise limb viability, potentially requiring further surgical intervention.
  \item \textbf{Cartilage Damage:} Iatrogenic injury to healthy articular cartilage during instrument insertion, manipulation, or inadvertent contact, which can accelerate degenerative changes.
  \item \textbf{Instrument Breakage:} Although rare, fragments of surgical instruments can break off within the joint, necessitating retrieval.
  \item \textbf{Synovial Fistula:} Persistent leakage of synovial fluid from a portal site, which may require prolonged dressing changes or, rarely, surgical closure.
  \item \textbf{Arthrofibrosis:} Excessive scar tissue formation within the joint, leading to stiffness, limited range of motion, and persistent pain, sometimes requiring further arthroscopic lysis of adhesions.
  \item \textbf{Failure of Repair/Reconstruction:} The repaired or reconstructed tissue may fail to heal, re-tear, or the graft may fail, leading to recurrence of symptoms and potentially requiring revision surgery.
  \item \textbf{Complex Regional Pain Syndrome (CRPS):} A rare but severe chronic pain condition characterized by disproportionate pain, swelling, and autonomic dysfunction in an extremity.
  \item \textbf{Compartment Syndrome:} A very rare but serious complication where increased pressure within a muscle compartment compromises blood flow, leading to tissue damage, requiring emergent fasciotomy.
  \item \textbf{Fluid Extravasation:} Leakage of irrigation fluid into surrounding soft tissues, which can cause swelling, nerve compression, or, rarely, airway compromise if it extends to the neck/chest.
\end{itemize}

\section*{Postoperative Care \& Recovery}
Immediately following arthroscopic repair, the patient is transferred to the Post-Anesthesia Care Unit (PACU) for close monitoring of vital signs, pain levels, and neurovascular status of the affected limb. Pain management is initiated, typically with a combination of oral or intravenous analgesics, and the surgical dressing is checked for excessive bleeding or swelling. Most arthroscopic procedures are performed on an outpatient basis, allowing patients to be discharged home the same day once they have fully recovered from anesthesia, their pain is adequately controlled, and they meet all discharge criteria.

During the first 24-48 hours at home, patients are generally advised to adhere strictly to the RICE protocol: Rest, Ice, Compression, and Elevation, to minimize swelling, bruising, and pain. Prescriptions for oral pain medication are provided, along with detailed instructions for wound care, including keeping the portal sites clean and dry. Depending on the specific procedure performed, weight-bearing restrictions (e.g., non-weight-bearing with crutches for knee procedures) or the use of slings or braces will be prescribed to protect the repair. Physical therapy often commences within the first week, focusing on gentle range-of-motion exercises and muscle activation to prevent stiffness and atrophy. Over the first 1-2 weeks, swelling and bruising gradually subside, and patients typically attend a follow-up appointment for wound check and suture or staple removal. Long-term recovery involves a structured and progressive physical therapy program, which can span from several weeks to many months, gradually increasing strength, flexibility, proprioception, and functional capacity, with a carefully guided return to activities as tolerated and approved by the surgeon and therapist. Full recovery and return to high-impact activities may take 6-12 months, particularly after ligament reconstructions.

During arthroscopic repair, the surgeon meticulously documents all intra-operative findings, which form a critical part of the patient's medical record. This includes a detailed description of the joint's condition, the specific pathology identified (e.g., type and location of meniscal tear, grade of cartilage damage, extent of rotator cuff tear), the presence of loose bodies, synovial inflammation, or any other abnormalities. The success of the repair, such as the stability of a meniscal repair, the tension of a ligament graft, or the security of a rotator cuff reattachment, is also noted. These findings directly inform the postoperative management and rehabilitation plan.

If tissue is resected during the procedure, such as a meniscal fragment, synovial tissue, or loose bodies, these specimens are sent to pathology for histological examination. The pathology report provides a definitive diagnosis, confirming the nature of the tissue (e.g., fibrocartilage for meniscus, inflamed synovium, osteochondral fragments). This report is crucial for understanding the underlying disease process, especially in cases of inflammatory arthropathy, pigmented villonodular synovitis, or suspected neoplastic conditions, although malignancy is rare in routine arthroscopic specimens. The correlation of surgical findings with the pathology report ensures a comprehensive understanding of the patient's condition and guides any further medical management.

A normal and successful postoperative course following arthroscopic repair is characterized by a predictable progression of healing and functional recovery. Initially, patients can expect mild to moderate pain, which is typically well-managed with prescribed oral analgesics and gradually diminishes over the first few days to weeks. Swelling and bruising around the joint are common but should gradually resolve with adherence to the RICE protocol. The small portal incisions heal quickly, usually within 10-14 days, with minimal scarring. Patients should experience a progressive improvement in joint range of motion and strength as they diligently follow their prescribed physical therapy program. The timeline for achieving specific milestones, such as full weight-bearing, independent ambulation, or return to sports, varies significantly depending on the specific procedure performed and individual patient factors, but generally follows a structured rehabilitation protocol. Ultimately, a successful course leads to a significant reduction or complete resolution of pre-operative symptoms, restoration of joint stability, and the ability to return to desired daily activities and, eventually, recreational or athletic pursuits without significant pain or limitation.

Postoperative care and diligent follow-up are crucial for optimizing the long-term success of arthroscopic repair and ensuring a smooth recovery.
\begin{itemize}
  \item \textbf{Wound Care:} Patients are instructed to keep the small portal incisions clean and dry, typically for the first 48-72 hours. Dressings are changed as directed, and signs of infection (redness, warmth, pus, increasing pain) are monitored.
  \item \textbf{Pain Management:} Oral analgesics are prescribed, and patients are advised on their proper use, often combined with ice and elevation.
  \item \textbf{Physical Therapy (PT):} A structured and progressive physical therapy program is almost universally prescribed, often commencing within the first week post-surgery. This is vital for restoring range of motion, strength, proprioception, and overall joint function. Adherence to PT protocols is a key determinant of outcome.
  \item \textbf{First Post-operative Visit:} Typically scheduled 10-14 days after surgery for a wound check, removal of sutures or staples, and initial assessment of swelling, pain, and early range of motion.
  \item \textbf{Subsequent Clinical Visits:} Follow-up appointments with the surgeon are usually scheduled at 1 month, 3 months, 6 months, and 1 year post-surgery to assess progress, evaluate joint stability, monitor functional recovery, and address any concerns.
  \item \textbf{Activity Resumption:} Gradual return to daily activities, work, and sports is carefully guided by the surgeon and physical therapist, following specific, procedure-dependent protocols to prevent re-injury. High-impact activities or competitive sports may be restricted for several months.
  \item \textbf{Warning Signs:} Patients are educated to contact their doctor immediately if they experience increasing or severe pain, spreading redness, warmth, pus drainage from the incision sites, fever (above 101.5$^{\circ}$F or 38.6$^{\circ}$C), persistent numbness or weakness, significant calf pain or swelling (suggesting DVT), or any sudden loss of joint function.
\end{itemize}
Adherence to these guidelines is essential for achieving the best possible functional outcome and minimizing complications.

\section*{Outcomes \& Prognosis}
Arthroscopic procedures are among the most frequently performed orthopedic surgeries globally, underscoring their broad utility and effectiveness in managing joint pathologies. Knee arthroscopy, in particular, is exceedingly common, with millions performed annually for conditions such as meniscal tears and anterior cruciate ligament (ACL) ruptures. The incidence of specific arthroscopic procedures varies significantly by age, sex, and activity level. For instance, ACL reconstruction is most prevalent in young, active individuals, especially athletes involved in pivoting sports, with a higher incidence observed in females for non-contact injuries. Rotator cuff repairs are more common in middle-aged and older adults, often associated with degenerative changes, repetitive overhead activities, or acute trauma. Meniscal tears can occur across all age groups, ranging from acute sports-related injuries in younger populations to degenerative tears in the elderly. While mortality rates directly associated with arthroscopic surgery are exceedingly low (estimated at <0.01%), morbidity, though generally minor, can include complications such as infection (reported rates typically 0.1-1%), deep vein thrombosis (0.5-2%), and nerve injury (0.1-0.5%), with rates varying depending on the specific joint, complexity of the procedure, and patient comorbidities.

The outcomes and prognosis following arthroscopic repair are generally favorable, with high success rates reported for many common procedures, leading to significant pain relief and functional improvement. For meniscal repair, success rates, defined by healing and absence of re-tear, can range from 70-90\%, influenced by tear type, location, patient age, and chronicity. ACL reconstruction typically boasts success rates of 85-95\% in restoring knee stability and enabling a return to pre-injury activity levels, including sports, although the risk of graft re-rupture exists, particularly in younger athletes. Rotator cuff repair outcomes are also largely positive, with good to excellent pain relief and functional improvement achieved in 80-90\% of patients, though re-tear rates, often asymptomatic, can vary.

Functional outcomes, including pain reduction, improved range of motion, enhanced joint stability, and the ability to return to desired activity levels, are the primary goals and are frequently achieved. Prognostic factors significantly influencing outcomes include the patient's age, overall health status, the severity and chronicity of the initial injury, the presence of concomitant arthritis or cartilage damage, and, critically, the patient's adherence to the prescribed postoperative rehabilitation program. While arthroscopy can effectively treat many specific conditions, it does not always prevent the progression of underlying degenerative joint disease in all cases, and recurrence of some pathologies, such as meniscal tears or labral tears, is possible depending on activity levels, biological healing capacity, and the presence of predisposing factors. Landmark studies, such as those evaluating outcomes of ACL reconstruction or rotator cuff repair, consistently demonstrate the efficacy of arthroscopic techniques in restoring joint function and improving quality of life.

\section*{References}
\begin{enumerate}
  \item MedlinePlus. Arthroscopy. U.S. National Library of Medicine; 2024. Available from: \url{https://medlineplus.gov/arthroscopy.html}
  \item Azar FM, Beaty JH, Canale ST. Campbell's Operative Orthopaedics. 14th ed. Elsevier; 2021.
  \item Miller MD, Thompson SR. DeLee and Drez's Orthopaedic Sports Medicine: Principles and Practice. 5th ed. Elsevier; 2020.
  \item American Academy of Orthopaedic Surgeons (AAOS). OrthoInfo. Available from: \url{https://orthoinfo.aaos.org/}
\end{enumerate}

\clearpage
